\documentclass[10pt,a4paper]{amsart}
\usepackage[margin=19mm]{geometry}
\usepackage{amsmath,amssymb,mathtools}
\usepackage[scr=rsfs,cal=euler]{mathalfa}
\usepackage{xfrac}
\usepackage[unicode,hidelinks,bookmarks=false]{hyperref}
\newcommand{\ie}{i.e.\ }
\newcommand{\cf}{cf.\ }
\newcommand{\resp}{respectively\ }
\newcommand{\Subsection}{\subsection}
\providecommand{\nobreakdash}{\nobreak\hbox{-}}
\setlength{\emergencystretch}{2em}
\multlinegap0pt
\usepackage[shortlabels]{enumitem}
\usepackage{mathtools}
\usepackage{amssymb}
\usepackage{float}
\usepackage{xcolor}
\newtheorem{proposition}{Proposition}
\title{Bia\l ynicki-Birula Decompositions of Nakajima Quiver Varieties, Quiver Chains and Star-Shaped Quivers}
\author{Juan Sebastian Numpaque-Roa}
\date{Source: arXiv:2607.17891v1 (CC BY 4.0)}
\begin{document}
\maketitle

\noindent\emph{Source and licence.} Bia\l ynicki-Birula Decompositions of Nakajima Quiver Varieties, Quiver Chains and Star-Shaped Quivers. Version arXiv:2607.17891v1 (CC BY 4.0), distributed by arXiv under the Creative Commons Attribution 4.0 International licence, http://creativecommons.org/licenses/by/4.0/, as recorded in the arXiv licence metadata for this exact version and shown on the abstract page https://arxiv.org/abs/2607.17891v1. Full licence notice: \texttt{licenses/arxiv-2607.17891-CC-BY-4.0.txt}.\par\smallskip

\noindent\textit{Editorial scope.} Counted unit: the article's proposition equivariantStructureUniversalFamily (source0.tex lines 616--619). The article's complete original proof of it is delivered with it (source0.tex lines 620--677, one \texttt{proof} environment of 58 lines). No mathematics is written, repaired or re-derived: the statement and the proof are the article's own lines, in the article's own notation, and no retained line is substituted or re-flowed.  Retained uncounted as the source-local context the proof needs: lines 535--537.  Nothing was omitted from any retained span, so the package carries no omission record.  No retained line contains a citation, so the wrapper carries no bibliography and no bibliography entry.  No retained line contains a figure environment, an image-inclusion command or drawing code, so no image is delivered and the package declares no asset dependency.  Editorial formatting only, in addition: an AMS \texttt{amsart} preamble that restates the article's own declarations and macros used by the retained lines, namely the environments \texttt{proposition} on separate counters, together with the standard packages those lines need.  The retained environments print in this document as Proposition 1 in order of appearance, whereas the article prints the same items with the numbering of its own full document.  The complete native source of the article as distributed in the arXiv e-print for this exact version is bundled unchanged as \texttt{sources/205607/source0.tex}.
	\begin{proposition}\label{existenceUniversalFamilyNakajimaQuiverVarieties}
		For indivisible dimension vector $\mathbf{d}$ and generic stability parameter $\theta$, $\mathcal{M}^{\theta}(\overline{Q},\mathbf{d})$ is a fine moduli space for families of stable Nakajima quiver representations and represents the functor $\mathcal{F}^{\theta}_{\mathbf{d}}$. 
	\end{proposition}

	\begin{proposition}\label{equivariantStructureUniversalFamily}
		There is a $\mathbb{C}^*$-equivariant structure on the universal family $\mathbb{M}^{\theta}(\overline{Q},\mathbf{d})$ acting with weight 0 on ordinary arrows
		$\alpha\in Q_1$ and weight one on opposite arrows $\overline\alpha$. 
	\end{proposition}

	\begin{proof}
		We keep the notation from Proposition \ref{existenceUniversalFamilyNakajimaQuiverVarieties}. Define a $\mathbb C^*$-action on each tautological vertex bundle $\widetilde{\mathcal{V}}_i$ by
		$$
		t\cdot(\mathbf{V},v):=(t\cdot \mathbf{V},v)
		$$
		and notice that this covers the $\mathbb C^*$-action on $\mu_{\mathbb{C}}^{-1}(0)^{\theta\text{-}s}$. We now see how this $\mathbb{C}^*$-action interacts with the tautological arrow maps. Let
		$\alpha$ be an ordinary arrow, then
		\begin{equation}\label{equivarianceArrowsBeforeDescending}
			\widetilde{\mathcal{X}}_\alpha(t\cdot(\mathbf{V},v))=(t\cdot\mathbf{V},X_{\alpha}(v))=t\cdot(\mathbf{V},X_{\alpha}(v))=t\cdot \widetilde{\mathcal{X}}_{\alpha}(\mathbf{V},v).
		\end{equation}
		
		For an opposite arrow $\overline{\alpha}$ we get
		\begin{equation}\label{equivarianceInverseArrowsBeforeDescending}
			\widetilde{\mathcal{Y}}_{\overline\alpha}(t\cdot(\mathbf{V},v))= (t\cdot \mathbf{V},tY_{\overline{\alpha}}(v))
			=t(t\cdot\widetilde{\mathcal{Y}}_{\overline{\alpha}}(\mathbf{V},v)).
		\end{equation}
		It remains to see that this \(\mathbb C^*\)-equivariant structure descends
		to the universal family $\mathbb{M}^{\theta}(\overline{Q},\mathbf{d})$. Recall, from Proposition \ref{existenceUniversalFamilyNakajimaQuiverVarieties}, that the $\text{GL}(\mathbf{d})$-action on $\widetilde{\mathcal{V}}_i$ is given by
		$$
		g\cdot(\mathbf{V},v)=(g\cdot \mathbf{V},\chi(g)g_i v)
		$$
		for some character $\chi$. Then
		$$
		t\cdot(g\cdot(\mathbf{V},v))=t\cdot(g\cdot \mathbf{V},\chi(g)g_i v)=
		(t\cdot(g\cdot \mathbf{V}),\chi(g)g_i v),
		$$
		while
		$$
		g\cdot(t\cdot(\mathbf{V},v))=
		g\cdot(t\cdot \mathbf{V},v)=
		(g\cdot(t\cdot \mathbf{V}),\chi(g)g_i v).
		$$
		Since the $\text{GL}(\mathbf{d})$-action on $\mu^{-1}_{\mathbb{C}}(0)^{\theta\text{-}s}$ commutes with the
		$\mathbb C^*$-action, this descends to the universal family. 
		Consequently, for every $t\in\mathbb{C}^*$ and every
		$[\mathbf{V}]\in \mathcal{M}^{\theta}(\overline{Q},\mathbf{d})$, we obtain isomorphisms
		$$
		\rho_i(t):
		(\mathbb{V}_i)_{\{[\mathbf{V}]\}}
		\to
		(\mathbb{V}_i)_{\{t\cdot[\mathbf{V}]\}}
		$$
		for all $i\in Q_0$. The relations in Equations (\ref{equivarianceArrowsBeforeDescending}) and (\ref{equivarianceInverseArrowsBeforeDescending}) now become respectively:
		$$
		\mathbb{X}_{\alpha}|_{t\cdot[\mathbf{V}]}
		\circ \rho_{t\alpha}(t)=\rho_{h\alpha}(t)\circ \mathbb{X}_{\alpha}|_{[\mathbf{V}]},
		$$
		and
		$$
		\mathbb{Y}_{\overline\alpha}|_{t\cdot[\mathbf{V}]}
		\circ \rho_{h\alpha}(t)=
		t\ 
		\rho_{t\alpha}(t)
		\circ \mathbb{Y}_{\overline\alpha}|_{[\mathbf{V}]}.
		$$
		Thus the universal family carries the desired
		\(\mathbb C^*\)-equivariant structure.
	\end{proof}

\end{document}
